\section{Evaluation Methodology}

\subsection{Benchmarks and provenance}

The initial study follows the AreCELearnedYet single-table benchmark lineage
and currently includes Census13, Forest10, Power7, and DMV11. Dataset adapters
validate source archive and content hashes, construct typed PostgreSQL tables,
and extract audited workloads. DMV11 uses authoritative external exact truth;
ordinary capture does not manufacture truth. The benchmark data root is
selected through \texttt{EXTSTATS\_RESEARCH\_DATA\_ROOT}, so an absolute local
path is not part of semantic run identity.

Every canonical RunManifest records the research repository and committed
research SHA, advisor repository and SHA, patched PostgreSQL repository and
SHA, stock PostgreSQL version or commit, dataset and workload identity, sample
rows and seed, statistics target, candidate limit, search budget, and relevant
experiment parameters. Canonical runs require a clean research working tree;
uncommitted experiment provenance is invalid.

\subsection{System roles and sample policy}

Stock PostgreSQL is the baseline and deployment target. The patched tree is
used only for design-time catalogless planner evaluation. A fixed typed sample
is created once per canonical run; native payloads are independently rebuilt on
the full database for transfer experiments. Ordinary statistics target,
sampling parameters, and relation population are held fixed within each
comparison.

\subsection{Metrics and comparisons}

RQ1 compares stock PostgreSQL 16, a high-statistics-target baseline, simple
selection heuristics, and \advisor. It reports arithmetic mean \,q-error,
p50, p95, p99, maximum, and improved/unchanged/worsened counts. Published
learned-CE results from \cite{Wang2021AreWeReady} are a clearly labeled
literature/contextual reference, not a head-to-head experiment under identical
systems and workloads. A gap-closure metric is reported only if the metric and
workload correspondence can be demonstrated.

RQ2 uses the existing full-data transfer protocol. P0 is a fresh full-data
baseline before deployment. P1 deploys the immutable Recommendation and runs
the final \texttt{ANALYZE}. P2 preserves the same ordinary-statistics state,
removes exactly the advisor-managed objects in a research-only transaction,
measures without \texttt{ANALYZE}, and rolls back. P1/P2 ordinary-statistics
fingerprints are expected to match; P0 can differ because deployment's final
\texttt{ANALYZE} refreshes the relation.

RQ3 compares hypothetical plans with physical deployment plans and relates
sample-sandbox utility to full-data deployment utility. RQ4 compares random-k,
workload-frequency top-k, dependency/correlation top-k, singleton-utility
top-k, greedy ADD, and exhaustive search on tiny candidate universes. RQ5
measures snapshot time and size, sample size, native construction, candidate
profiling, search time and planner evaluations, deployment DDL/\texttt{ANALYZE},
storage, planning-time change, and refresh cost. We do not claim zero online
overhead; the narrower statement is that no separate learned-model inference
path is added to production query processing.

\subsection{Preliminary versus confirmatory evidence}

Existing Census13, Forest10, Power7, and DMV11 runs are pilot/preliminary
evidence that the end-to-end pipeline and full-data validation path work. The
paper will only promote a result to a final claim when its immutable raw and
derived artifacts, manifest, and source identities pass the repository audit.
The current results section therefore defines schemas and TODOs rather than
inventing values.

\subsection{Reproducibility}

The research harness stores compact derived evidence while referring to large
immutable bundles by run ID and digest. Reproduction starts from a clean
committed research SHA, verifies advisor and PostgreSQL identities, reconstructs
the frozen input artifacts, and reruns the declared experiment command. The
paper's companion experiment plan is the specification for the missing
confirmatory runs.
